% topics/curvilinear_coordinates/main.tex
% Copyright (C) 2026 Hirokazu Maesaka (RIKEN SPring-8 Center)
% SPDX-License-Identifier: CC-BY-4.0
%
% Bilingual memo. Build with the commands in AGENTS.md (English: \memolang=en, Japanese: \memolang=ja).

% common/tex/preamble.tex
% Copyright (C) 2026 Hirokazu Maesaka (RIKEN SPring-8 Center)
% SPDX-License-Identifier: CC-BY-4.0
%
% Shared preamble for all topic documents. Loaded by topics/<topic>/main.tex with \input.
% Build with lualatex; the language is selected with \memolang (en or ja) through latexmk -pretex.

\documentclass[11pt,a4paper]{article}

% Language switch. The default is English. latexmk -pretex='\def\memolang{ja}' overrides it.
\providecommand{\memolang}{en}
% \ifx compares macro meanings. Unlike \ifdefstring, it is expandable, so \en and \ja also work inside \edef.
\def\memolangEN{en}
\def\memolangJA{ja}
\newcommand{\en}[1]{\ifx\memolang\memolangEN #1\fi}
\newcommand{\ja}[1]{\ifx\memolang\memolangJA #1\fi}

\usepackage{luatexja-preset}
\usepackage{etoolbox}
\usepackage{amsmath,amssymb}
\usepackage[style=numeric,backend=biber,sorting=none]{biblatex}
\addbibresource{common/references/references.bib}

\setlength{\parindent}{0pt}
\setlength{\parskip}{0.6em}

% Slide markers (write build/<topic>/<jobname>.slides for common/slides/build_slides.py).
% common/tex/slides.tex
% Copyright (C) 2026 Hirokazu Maesaka (RIKEN SPring-8 Center)
% SPDX-License-Identifier: CC-BY-4.0
%
% Slide markers. Each marker writes one line "kind:payload" to build/<topic>/<jobname>.slides during compilation.
% common/slides/build_slides.py reads that file and builds the PowerPoint deck. The markers print nothing in the PDF.
% Arguments are expanded at the point of use, so \en{...} and \ja{...} select the language of the current build.
% Do not use fragile commands (for example \cite or \label) inside slide arguments.

\newwrite\slidesfile
\AtBeginDocument{\immediate\openout\slidesfile=\jobname.slides\relax}
\AtEndDocument{\immediate\closeout\slidesfile}

% #1: entry kind, #2: payload. The payload is expanded, then written as plain characters.
\newcommand{\slideentry}[2]{%
  \begingroup
  \edef\slidetmp{#2}%
  \edef\slidedet{\detokenize\expandafter{\slidetmp}}%
  \immediate\write\slidesfile{#1:\slidedet}%
  \endgroup}

% Raw entries are written without expansion. Use them for math and paths, which contain fragile commands.
\newcommand{\slideentryraw}[2]{%
  \immediate\write\slidesfile{#1:\detokenize{#2}}}

\newcommand{\slidetitle}[1]{\slideentry{title}{#1}}
\newcommand{\slidepoint}[1]{\slideentry{point}{#1}}
\newcommand{\slidemath}[1]{\slideentryraw{math}{#1}}
\newcommand{\slidefigure}[1]{\slideentryraw{figure}{#1}}
\newcommand{\slidenote}[1]{\slideentry{note}{#1}}



\title{\en{Curvilinear (Frenet-Serret) coordinates and the Hamiltonian of a bending magnet}%
       \ja{曲線座標（Frenet-Serret 座標）における偏向磁石のハミルトニアン}}
\author{Hirokazu Maesaka}
\date{2026}

\begin{document}
\maketitle

% Slide deck source. These markers print nothing in the PDF; common/slides/build_slides.py builds the deck from them.
\slidetitle{\en{Curvilinear coordinates and the Hamiltonian of a bending magnet}%
            \ja{曲線座標における偏向磁石のハミルトニアン}}
\slidepoint{\en{Frenet-Serret frame, s as the independent variable}%
            \ja{Frenet-Serret 座標系と，s を独立変数とする定式化}}
\slidenote{\en{Derivation follows [lee2019, Chap. 2] and [wiedemann2015, p. 93].}%
           \ja{導出は [lee2019, 第 2 章] と [wiedemann2015, p. 93] に従う。}}

\slidetitle{\en{Frenet-Serret frame and metric}\ja{Frenet-Serret 座標系と計量}}
\slidepoint{\en{Curvature h = 1/ρ; x outward, y vertical}%
            \ja{曲率 h = 1/ρ。x は外向き，y は鉛直}}
\slidemath{\mathrm{d}\ell^2 = (1+hx)^2\,\mathrm{d}s^2 + \mathrm{d}x^2 + \mathrm{d}y^2}
\slidenote{\en{The factor 1+hx is the only non-trivial metric coefficient.}%
           \ja{非自明な計量係数は 1+hx だけである。}}

\slidetitle{\en{Hamiltonian with s as the independent variable}\ja{s を独立変数とするハミルトニアン}}
\slidepoint{\en{Legendre transform of the Lagrangian with s as the independent variable}%
            \ja{s を独立変数とするラグランジアンの Legendre 変換}}
\slidemath{K = -(1+hx)\Big[\sqrt{\tfrac{E^2}{c^2} - m^2c^2 - P_x^2 - P_y^2} + qA_s\Big]}

\slidetitle{\en{Normalized Hamiltonian}\ja{規格化されたハミルトニアン}}
\slidepoint{\en{Normalized with P0; delta is a constant of motion for a static B-field}%
            \ja{P0 で規格化。静磁場では delta は運動の定数}}
\slidemath{H = -(1+hx)\sqrt{(1+\delta)^2 - (p_x-a_x)^2 - (p_y-a_y)^2} - (1+hx)\,a_s}

\slidetitle{\en{Ideal dipole potential}\ja{理想的な偏向磁石のポテンシャル}}
\slidepoint{\en{Uniform vertical field B0; the sign is fixed by the reference-orbit equilibrium}%
            \ja{一様な鉛直磁場 B0。符号は基準軌道の平衡条件で決まる}}
\slidemath{-(1+hx)\,a_s = h\big(x + \tfrac{1}{2} h x^2\big)}

\slidetitle{\en{Equations of motion}\ja{運動方程式}}
\slidemath{x' = \frac{(1+hx)\,p_x}{\sqrt{(1+\delta)^2 - p_x^2 - p_y^2}}}
\slidemath{p_x' = -h(1+hx) + h\sqrt{(1+\delta)^2 - p_x^2 - p_y^2} - kx}
\slidenote{\en{Hamilton's equations with s as the independent variable.}%
           \ja{s を独立変数とする正準方程式である。}}

\slidetitle{\en{Linearized equations of motion}\ja{線形化された運動方程式}}
\slidemath{x'' + (h^2 + k)\,x = h\,\delta}
\slidemath{y'' - k\,y = 0}
\slidepoint{\en{The bend adds the inhomogeneous term h delta, which generates the dispersion}%
            \ja{偏向により非斉次項 h delta が加わり，分散が生じる}}

\slidetitle{\en{Checks and open points}\ja{確認事項と未解決の点}}
\slidepoint{\en{Derivations checked with sympy in the notebook}\ja{導出はノートブックで sympy により検証済み}}
\slidepoint{\en{RF and quadrupole curvature corrections are not included}\ja{RF と四極磁石の曲率補正は含まない}}
\slidepoint{\en{Check the tau canonical transformation against [lee2019, Chap. 2]}%
            \ja{tau の正準変換は [lee2019, 第 2 章] で確認が必要}}

\section{\en{Scope and conventions}\ja{対象と記号の規約}}

\en{This memo derives the Hamiltonian in a curvilinear frame that follows the ideal (reference) orbit through a
planar bending magnet, and then the equations of motion (EOM) around the reference orbit. Results are checked with
sympy in the companion notebook \texttt{curvilinear\_coordinates.py}. The derivation follows
\cite[Chap.~2]{lee2019} and \cite[p.~93]{wiedemann2015}. The page-level content of those references has not been
re-checked against this memo.}%
\ja{本メモでは，理想軌道（基準軌道）に沿って偏向磁石を通る曲線座標系でハミルトニアンを導き，基準軌道の周りの運動方程式を求める。
結果は付随するノートブック \texttt{curvilinear\_coordinates.py} で sympy を用いて検証している。導出は
\cite[第 2 章]{lee2019} と \cite[p.~93]{wiedemann2015} を参考にしている。
ただし，それらの文献の該当ページの内容は本メモ執筆時に再確認していない。}

\en{Conventions. $s$ is the arc length along the reference orbit and the independent variable. $x$ is the
horizontal coordinate, positive outward. $y$ is the vertical coordinate. $h(s)=1/\rho(s)$ is the curvature of the
reference orbit, with $h>0$ when the center of curvature lies inside. $q$ and $m$ are the charge and the mass. $P_0$
and $\beta_0$ are the reference momentum and the reference velocity factor. Some texts use other signs for $x$ and
$h$; check the sign convention before comparing formulas.}%
\ja{記号の規約。$s$ は基準軌道に沿った弧長で，独立変数である。$x$ は外向きを正とする水平座標，$y$ は鉛直座標である。
$h(s)=1/\rho(s)$ は基準軌道の曲率で，曲率中心が内側にあるとき $h>0$ とする。$q$ と $m$ は電荷と質量，$P_0$ と $\beta_0$ は
基準運動量と基準速度係数である。文献によって $x$ と $h$ の符号規約が異なるため，式を比較する際は符号規約を確認すること。}

\section{\en{Frenet-Serret frame and metric}\ja{Frenet-Serret 座標系と計量}}

\en{For a planar curve $\mathbf{R}(s)$ with curvature $h(s)$ and zero torsion, the Frenet-Serret equations in this
convention are}%
\ja{曲率 $h(s)$，ねじれゼロの平面曲線 $\mathbf{R}(s)$ に対して，本メモの規約での Frenet-Serret の式は次のとおりである。}
\begin{align}
  \frac{d\mathbf{T}}{ds} &= -h\,\mathbf{e}_x, &
  \frac{d\mathbf{e}_x}{ds} &= h\,\mathbf{T}, &
  \mathbf{T} &= \frac{d\mathbf{R}}{ds},
\end{align}
\en{where $\mathbf{T} = d\mathbf{R}/ds$ is the unit tangent vector of $\mathbf{R}(s)$, $\mathbf{e}_x$ is the outward
horizontal unit normal, and $\mathbf{e}_y$ is the vertical unit vector. A particle near the reference orbit is at}%
\ja{ここで $\mathbf{T} = d\mathbf{R}/ds$ は $\mathbf{R}(s)$ の単位接線ベクトル，$\mathbf{e}_x$ は外向きの水平単位法線，
$\mathbf{e}_y$ は鉛直単位ベクトルである。基準軌道の近傍の粒子の位置は}
\begin{equation}
  \mathbf{r}(s, x, y) = \mathbf{R}(s) + x\,\mathbf{e}_x(s) + y\,\mathbf{e}_y .
\end{equation}
\en{The line element is}\ja{線要素は}
\begin{equation}
  d\ell^2 = (1 + hx)^2\,ds^2 + dx^2 + dy^2 .
  \label{eq:metric}
\end{equation}
\en{The factor $1+hx$ is the only non-trivial metric coefficient. The Hamiltonian is therefore not a Cartesian one.
The metric is checked for a circular orbit in the notebook.}%
\ja{非自明な計量係数は $1+hx$ だけである。したがってハミルトニアンはデカルト座標の形にはならない。
円軌道に対する計量は，ノートブックで確認している。}

\section{\en{Lagrangian and Hamiltonian with $s$ as the independent variable}%
         \ja{$s$ を独立変数とするラグランジアンとハミルトニアン}}

\en{With time $t$ as the independent variable, the Lagrangian of a particle in a static magnetic field is}%
\ja{時間 $t$ を独立変数とすると，静磁場中の粒子のラグランジアンは}
\begin{equation}
  L = -mc^2\sqrt{1 - \frac{v^2}{c^2}} + q\,\mathbf{v}\cdot\mathbf{A} .
\end{equation}
\en{The Euler-Lagrange equation for the Cartesian position $\mathbf{r}$ is}%
\ja{デカルト座標の位置 $\mathbf{r}$ に対するオイラー・ラグランジュ方程式は}
\begin{equation}
  \frac{d}{dt}\frac{\partial L}{\partial \mathbf{v}} - \frac{\partial L}{\partial \mathbf{r}} = 0,
  \qquad \mathbf{v} = \dot{\mathbf{r}} .
  \label{eq:el}
\end{equation}
\en{The velocity derivative is}%
\ja{速度による微分は}
\begin{equation}
  \frac{\partial L}{\partial \mathbf{v}} = \gamma m\,\mathbf{v} + q\,\mathbf{A},
  \qquad \gamma = \frac{1}{\sqrt{1 - v^2/c^2}} ,
  \label{eq:el_velocity}
\end{equation}
\en{and the position derivative, with $\mathbf{v}$ held fixed, is}%
\ja{であり，$\mathbf{v}$ を固定した位置による微分は}
\begin{equation}
  \frac{\partial L}{\partial \mathbf{r}} = q\,\nabla(\mathbf{v}\cdot\mathbf{A})
  = q\big[(\mathbf{v}\cdot\nabla)\mathbf{A} + \mathbf{v}\times(\nabla\times\mathbf{A})\big] .
  \label{eq:el_position}
\end{equation}
\en{The vector potential is evaluated along the trajectory, so its total time derivative is}%
\ja{ベクトルポテンシャルは軌道上で評価するので，その全時間微分は}
\begin{equation}
  \frac{d\mathbf{A}}{dt} = \frac{\partial \mathbf{A}}{\partial t} + (\mathbf{v}\cdot\nabla)\mathbf{A} .
  \label{eq:dA_dt}
\end{equation}
\en{Substituting \eqref{eq:el_velocity}--\eqref{eq:dA_dt} into \eqref{eq:el}, the directional derivative
$(\mathbf{v}\cdot\nabla)\mathbf{A}$ of the vector potential along the trajectory cancels and}%
\ja{\eqref{eq:el_velocity}--\eqref{eq:dA_dt} を \eqref{eq:el} に代入すると，軌道に沿ったベクトルポテンシャルの
方向微分 $(\mathbf{v}\cdot\nabla)\mathbf{A}$ が相殺して}
\begin{equation}
  \frac{d}{dt}(\gamma m\,\mathbf{v}) = q\,\mathbf{v}\times\mathbf{B} - q\,\frac{\partial \mathbf{A}}{\partial t},
  \qquad \mathbf{B} = \nabla\times\mathbf{A} .
  \label{eq:lorentz}
\end{equation}
\en{For a static magnetic field $\partial\mathbf{A}/\partial t = 0$, and \eqref{eq:lorentz} is the Lorentz force
equation $d\mathbf{p}/dt = q\,\mathbf{v}\times\mathbf{B}$ with the mechanical momentum
$\mathbf{p} = \gamma m\mathbf{v}$. The same Euler-Lagrange procedure applied to $L_s$ with $s$ as the independent
variable gives the equations of motion used in the following sections. The notebook verifies the Cartesian
identity \eqref{eq:lorentz} with sympy.}%
\ja{静磁場では $\partial\mathbf{A}/\partial t = 0$ であるから，\eqref{eq:lorentz} は力学的運動量
$\mathbf{p} = \gamma m\mathbf{v}$ を用いてローレンツ力の運動方程式
$d\mathbf{p}/dt = q\,\mathbf{v}\times\mathbf{B}$ にほかならない。$s$ を独立変数とする $L_s$ に同じ手続きを
適用すると，以降の節で用いる運動方程式が得られる。デカルト座標での恒等式 \eqref{eq:lorentz} は，
ノートブックで sympy により検証している。}

\en{The action does not depend on the choice of the independent variable, so $L\,dt = L_s\,ds$. With
$t' = dt/ds$, $x' = dx/ds$, and $y' = dy/ds$, the velocity components satisfy $v_T t' = 1+hx$, $v_x t' = x'$, and
$v_y t' = y'$. Using these and the metric \eqref{eq:metric}, the Lagrangian with $s$ as the independent variable is}%
\ja{作用は独立変数の取り方によらないので $L\,dt = L_s\,ds$ である。$t' = dt/ds$，$x' = dx/ds$，$y' = dy/ds$ とすると，
速度成分は $v_T t' = 1+hx$，$v_x t' = x'$，$v_y t' = y'$ を満たす。これらと計量 \eqref{eq:metric} を用いると，
$s$ を独立変数とするラグランジアンは}
\begin{equation}
  L_s = -mc^2\sqrt{t'^2 - \frac{(1+hx)^2 + x'^2 + y'^2}{c^2}}
        + q\big[(1+hx)A_s + x'A_x + y'A_y\big] .
\end{equation}
\en{The variable $t'$ is not independent. The tangential velocity is $v_T = (1+hx)\,ds/dt = (1+hx)/t'$, and the
tangential mechanical momentum is $P_T = \gamma m v_T = E v_T/c^2$, so}%
\ja{変数 $t'$ は独立ではない。接線方向の速度は $v_T = (1+hx)\,ds/dt = (1+hx)/t'$ であり，接線方向の力学的運動量は
$P_T = \gamma m v_T = E v_T/c^2$ であるから}
\begin{equation}
  t' = \frac{(1+hx)\,E}{P_T c^2},
  \qquad P_T = \sqrt{\frac{E^2}{c^2} - m^2c^2 - P_x^2 - P_y^2} .
  \label{eq:tprime}
\end{equation}
\en{The generalized coordinates are $(t, x, y)$, and their conjugate momenta are}%
\ja{一般化座標は $(t, x, y)$ であり，その共役運動量は}
\begin{equation}
  \begin{aligned}
    p_t &= \frac{\partial L_s}{\partial t'} = -\gamma m c^2 = -E, \\
    p_x &= \frac{\partial L_s}{\partial x'} = \gamma m v_x + qA_x, \qquad
    p_y = \frac{\partial L_s}{\partial y'} = \gamma m v_y + qA_y .
  \end{aligned}
  \label{eq:momenta}
\end{equation}
\en{The Hamiltonian with $s$ as the time-like variable is $K = p_x x' + p_y y' + p_t t' - L_s$. Writing the mechanical
transverse momenta as $P_x = p_x - qA_x$ and $P_y = p_y - qA_y$, the Legendre transform gives}%
\ja{$s$ を時間的変数とするハミルトニアンは $K = p_x x' + p_y y' + p_t t' - L_s$ である。力学的な横方向運動量を
$P_x = p_x - qA_x$，$P_y = p_y - qA_y$ とすると，Legendre 変換により次の式が得られる。}
\begin{equation}
  K = -(1+hx)\Big[\sqrt{\frac{E^2}{c^2} - m^2c^2 - P_x^2 - P_y^2} + qA_s\Big].
  \label{eq:K}
\end{equation}
\en{The prefactor $-(1+hx)$ comes from $t'$ and the metric factor.}%
\ja{前因子 $-(1+hx)$ は $t'$ と計量因子に由来する。}

\section{\en{Normalization}\ja{規格化}}

\en{Normalize with the reference momentum $P_0$: $1+\delta = p/P_0$, $p_x \to P_x/P_0$, $p_y \to P_y/P_0$, and
$a_x = qA_x/P_0$, $a_y = qA_y/P_0$, $a_s = qA_s/P_0$. Dividing \eqref{eq:K} by $P_0$ gives}%
\ja{基準運動量 $P_0$ で規格化する。すなわち $1+\delta = p/P_0$，$p_x \to P_x/P_0$，$p_y \to P_y/P_0$，
$a_x = qA_x/P_0$，$a_y = qA_y/P_0$，$a_s = qA_s/P_0$ とする。\eqref{eq:K} を $P_0$ で割ると}
\begin{equation}
  H = -(1+hx)\sqrt{(1+\delta)^2 - (p_x - a_x)^2 - (p_y - a_y)^2} \;-\; (1+hx)\,a_s .
  \label{eq:H_general}
\end{equation}
\en{The time-like variable $\tau$ is conjugate to $\delta$. It comes from $(t,E)$ by a canonical transformation
(the standard choice in \cite{lee2019} and \cite{wiedemann2015}). For a static magnetic field $H$ does not depend
on $\tau$, so $\delta$ is a constant of motion.}%
\ja{時間的変数 $\tau$ は $\delta$ と共役である。$(t,E)$ からの正準変換により得られる
（\cite{lee2019} と \cite{wiedemann2015} の標準的な選択）。静磁場では $H$ は $\tau$ に依らないため，$\delta$ は運動の定数である。}

\en{If the energy deviation $\delta_L = \Delta E/(P_0c)$ is used instead, the momentum ratio becomes}%
\ja{エネルギー偏差 $\delta_L = \Delta E/(P_0c)$ を用いる場合，運動量比は次のとおりである。}
\begin{equation}
  \Big(\frac{p}{P_0}\Big)^2 = 1 + \frac{2\,\delta_L}{\beta_0} + \delta_L^2 ,
\end{equation}
\en{which reduces to $(1+\delta)^2$ at first order.}%
\ja{これは 1 次の近似で $(1+\delta)^2$ に一致する。}

\section{\en{Vector potentials of the dipole and the quadrupole}%
         \ja{偏向磁石と四極磁石のベクトルポテンシャル}}

\en{An ideal sector dipole has a uniform vertical field $B_y = B_0$. With $A_x = 0$, the vertical curl in the
curvilinear frame gives $(1+hx)\,a_s = -h\,(x + \tfrac{1}{2} h x^2)$, that is}%
\ja{理想的な扇型偏向磁石は一様な鉛直磁場 $B_y = B_0$ を持つ。$A_x = 0$ とすると，曲線座標での鉛直方向の回転から
$(1+hx)\,a_s = -h\,(x + \tfrac{1}{2} h x^2)$ が得られ，すなわち}
\begin{equation}
  a_s = -\frac{h\,(x + \tfrac{1}{2} h x^2)}{1+hx}, \qquad h = \frac{qB_0}{P_0} = \frac{1}{\rho}.
  \label{eq:as_dip}
\end{equation}
\en{The overall sign is a convention. It is fixed by requiring the reference orbit to be an equilibrium of $H$.
With \eqref{eq:as_dip} the dipole term is exact:}%
\ja{符号は規約である。基準軌道が $H$ の平衡点であることを要求して決める。\eqref{eq:as_dip} により，偏向磁石の項は厳密に次のとおりである。}
\begin{equation}
  -(1+hx)\,a_s = h\,\big(x + \tfrac{1}{2} h x^2\big).
\end{equation}
\en{For a straight quadrupole with normalized gradient $k(s) = qg/P_0$ (focusing in $x$ for $k>0$), the potential
term needed for linear optics is $\tfrac{1}{2}k\,(x^2 - y^2)$. Curvature corrections to the quadrupole term are
higher order and are not kept.}%
\ja{規格化された勾配 $k(s) = qg/P_0$ を持つ直線四極磁石（$k>0$ で $x$ 方向に集束）では，線形光学に必要なポテンシャル項は
$\tfrac{1}{2}k\,(x^2 - y^2)$ である。四極項の曲率補正は高次であり，本メモでは扱わない。}

\section{\en{Full Hamiltonian and exact equations of motion}\ja{完全なハミルトニアンと厳密な運動方程式}}

\en{Collecting the terms, the Hamiltonian is}\ja{項をまとめると，ハミルトニアンは次のとおりである。}
\begin{equation}
  H(x,p_x,y,p_y,\delta;s) = -(1+hx)\sqrt{(1+\delta)^2 - p_x^2 - p_y^2}
    + h\big(x + \tfrac{1}{2} h x^2\big) + \tfrac{1}{2} k\,(x^2 - y^2).
  \label{eq:H_full}
\end{equation}
\en{Hamilton's equations with $s$ as the independent variable are}%
\ja{$s$ を独立変数とする正準方程式は次のとおりである。}
\begin{equation}
  x' = \frac{\partial H}{\partial p_x}, \quad
  p_x' = -\frac{\partial H}{\partial x}, \quad
  y' = \frac{\partial H}{\partial p_y}, \quad
  p_y' = -\frac{\partial H}{\partial y}, \quad
  \delta' = 0 .
\end{equation}
\en{The reference orbit is an equilibrium because $\partial H/\partial x = 0$ at $x = p = \delta = 0$. Explicitly,}%
\ja{基準軌道は $x = p = \delta = 0$ で $\partial H/\partial x = 0$ となるため平衡点である。具体的には，}
\begin{align}
  x' &= \frac{(1+hx)\,p_x}{\sqrt{(1+\delta)^2 - p_x^2 - p_y^2}}, \\
  p_x' &= -h(1+hx) + h\sqrt{(1+\delta)^2 - p_x^2 - p_y^2} - kx .
\end{align}
\en{and $y' = (1+hx)p_y/\sqrt{\cdots}$, $p_y' = ky$.}%
\ja{であり，$y' = (1+hx)p_y/\sqrt{\cdots}$，$p_y' = ky$ である。}

\section{\en{Linearized equations of motion}\ja{線形化された運動方程式}}

\en{Let $\mathbf{z} = (x, p_x, y, p_y, \delta)^\mathsf{T}$. The Jacobian of the exact EOM at the origin gives}%
\ja{$\mathbf{z} = (x, p_x, y, p_y, \delta)^\mathsf{T}$ とする。原点での厳密な運動方程式のヤコビアンは}
\begin{equation}
  \mathbf{z}' = \mathbf{A}\,\mathbf{z}, \qquad
  \mathbf{A} = \begin{pmatrix}
    0 & 1 & 0 & 0 & 0 \\
    -(h^2+k) & 0 & 0 & 0 & h \\
    0 & 0 & 0 & 1 & 0 \\
    0 & 0 & k & 0 & 0 \\
    0 & 0 & 0 & 0 & 0
  \end{pmatrix}.
\end{equation}
\en{The corresponding equations are}\ja{対応する方程式は次のとおりである。}
\begin{equation}
  x'' + (h^2 + k)\,x = h\,\delta, \qquad y'' - k\,y = 0 .
  \label{eq:linear}
\end{equation}
\en{The first equation is the Hill equation with the inhomogeneous term $h\delta$ from the bend. The second has no
$h^2$ term, as expected for a horizontal bend.}%
\ja{第 1 式は偏向による非斉次項 $h\delta$ を持つ Hill 方程式である。第 2 式は水平偏向のため $h^2$ 項を持たない。}

\en{Checks.}\ja{確認。}
\begin{itemize}
  \item \en{For constant $h$ and $k$, the constant $D = h/(h^2+k)$ satisfies $D'' + (h^2+k)D = h$, the uniform
    dispersion of a bending section.}%
    \ja{一定の $h$ と $k$ では，定数 $D = h/(h^2+k)$ が $D'' + (h^2+k)D = h$ を満たす。これは一様な偏向部の分散である。}
  \item \en{There is no linear coupling in this Hamiltonian, because there are no skew terms.}%
    \ja{歪み項がないため，このハミルトニアンに線形結合はない。}
\end{itemize}

\section{\en{Summary and open points}\ja{まとめと未解決の点}}

\begin{enumerate}
  \item \en{Metric: $d\ell^2 = (1+hx)^2 ds^2 + dx^2 + dy^2$.}%
    \ja{計量: $d\ell^2 = (1+hx)^2 ds^2 + dx^2 + dy^2$。}
  \item \en{Hamiltonian: $H = -(1+hx)\sqrt{(1+\delta)^2-(p_x-a_x)^2-(p_y-a_y)^2} - (1+hx)a_s$.}%
    \ja{ハミルトニアン: $H = -(1+hx)\sqrt{(1+\delta)^2-(p_x-a_x)^2-(p_y-a_y)^2} - (1+hx)a_s$。}
  \item \en{Linear EOM: $x'' + (h^2+k)x = h\delta$ and $y'' - ky = 0$, with $\delta' = 0$ in the absence of RF.}%
    \ja{線形方程式: $x'' + (h^2+k)x = h\delta$，$y'' - ky = 0$。RF がない場合 $\delta' = 0$。}
  \item \en{The RF cavity adds an electric potential, so $\delta$ and $\tau$ are no longer constant. Not covered here.}%
    \ja{RF 空洞は電位を加えるため，$\delta$ と $\tau$ は定数ではなくなる。本メモでは扱わない。}
  \item \en{Quadrupole curvature corrections and sextupole and octupole terms are not included.}%
    \ja{四極磁石の曲率補正，六極・八極の項は含めていない。}
  \item \en{Whether the $\tau$-canonical transformation gives exactly the form above should be checked against
    \cite[Chap.~2]{lee2019} before using $H$ beyond the linear level.}%
    \ja{$H$ を線形の範囲を超えて用いる前に，$\tau$ の正準変換により上記の形が厳密に得られるかを
    \cite[第 2 章]{lee2019} で確認すること。}
\end{enumerate}

\printbibliography[title={\en{References}\ja{参考文献}}]

\section*{\en{License}\ja{ライセンス}}
\en{This document is licensed under CC BY 4.0. Code snippets are GPL-3.0-or-later. See \texttt{LICENSING.md}.}%
\ja{本文書は CC BY 4.0 で提供する。コードの断片は GPL-3.0-or-later である。詳細は \texttt{LICENSING.md} を参照すること。}

\end{document}
